% ══════════════════════════════════════════════════════════════════
%  {{CARPETA}} — teoría
%  Módulo SIN preámbulo: se incluye desde main.tex con \include{}.
% ══════════════════════════════════════════════════════════════════

\portadaclase{{{TEMA_TITULO}}}{{{NN}}}{[N]}{[Nombre de la unidad]}

\section{{{TEMA_TITULO}}}

% ── Motivación ─────────────────────────────────────────────────────
% Relevancia del tema en exámenes de admisión (UNI, San Marcos, PUCP…)

% ── 1. Definiciones ────────────────────────────────────────────────
\begin{cajadef}[: concepto principal]
  % Definición formal del concepto.
\end{cajadef}

% ── 2. Propiedades ─────────────────────────────────────────────────
\begin{cajaprop}[1]
  % Propiedad o teorema con su condición de validez.
\end{cajaprop}

% ── 3. Fórmulas importantes ────────────────────────────────────────
\begin{cajaformula}
  \[ % fórmula clave \]
\end{cajaformula}

% ── 4. Casos especiales y observaciones ────────────────────────────

% ── 5. Errores comunes ─────────────────────────────────────────────
\begin{cajaadvert}
  % Error frecuente de los postulantes y cómo evitarlo.
\end{cajaadvert}

% ── 6. Ejemplos resueltos (básico, intermedio, admisión) ───────────
\begin{cajaejemplo}[1 — nivel básico]
  % Enunciado y solución paso a paso.
\end{cajaejemplo}

% ── 7. Resumen + tip final ─────────────────────────────────────────
\begin{cajatip}
  % Estrategia práctica para el examen.
\end{cajatip}
